\section{Implementation}
\label{sec:implementation}

ICD and AICD are implemented in the open-source BNNR
library~\citep{walo2026bnnr} (MIT licensed, available on PyPI as
\texttt{bnnr}).  The saliency back end is the
\texttt{pytorch-grad-cam} package~\citep{gildenblat2021gradcam}, so the maps of
Eq.~\eqref{eq:gradcam} come from a standard, widely used GradCAM implementation
rather than a bespoke one; by default the library enables that package's
eigen-smoothing option, as noted in Section~\ref{sec:method:saliency}.  This
section describes how the methods are applied during training.

\subsection{Online application and the saliency cache}
\label{sec:implementation:cache}

Because the saliency map depends on the current parameters, a naive
implementation would run a GradCAM pass for every image at every step, roughly
doubling the cost of training.  BNNR amortises this with a disk-backed cache.
Before a training run (or after a warm-up phase, as discussed in
Section~\ref{sec:discussion:limits}), the saliency map for each training image
is computed once and stored, keyed by a hash of the image content:
\begin{verbatim}
    from bnnr.xai_cache import XAICache
    cache = XAICache("./xai_cache")
    cache.precompute_cache(
        model=model,
        train_loader=train_loader,   # yields (image, label, index)
        target_layers=target_layers, # e.g. [model.layer4[-1]]
        n_samples=len(train_dataset),
        method="gradcam",
    )
\end{verbatim}
During training the augmentations read tile masks from the cache instead of
recomputing GradCAM, reducing the per-step overhead to a cache lookup and the
masking arithmetic of Eqs.~\eqref{eq:tilescore}--\eqref{eq:fill}.  The cache
reflects the model at the time of precomputation; refreshing it periodically
trades additional compute for saliency that tracks the evolving model more
closely.

\subsection{Use in a training loop}
\label{sec:implementation:loop}

The augmentations are plain callables that take a batch and its labels.  A
classification step looks as follows:
\begin{verbatim}
    from bnnr.icd import ICD, AICD
    icd  = ICD(model=model,  target_layers=target_layers,
               cache=cache,  threshold_percentile=70.0,
               tile_size=8,  fill_strategy="gaussian_blur",
               probability=0.5)
    aicd = AICD(model=model, target_layers=target_layers,
                cache=cache, threshold_percentile=70.0)

    for images, labels, index in train_loader:   # images uint8 NHWC
        images = icd.apply_batch_with_labels(
            images, labels, sample_indices=index)
        # ... convert to [0,1] BCHW, forward, backward, step ...
\end{verbatim}
The label passed to the augmentation is the class whose saliency defines the
mask; by default this is the ground-truth label, though the model's prediction
may be used instead.  The \texttt{DataLoader} yields a sample index as a third
element so the cache can retrieve the correct precomputed map.  Images are
expected in $[0,1]$ (or uint8, converted internally); ImageNet normalisation, if
used, is applied after the augmentation so that the saliency back end and the
fill operate on unnormalised pixels.

ICD and AICD are independent augmentations and need not be used together.  Each
can be inserted into any existing PyTorch loop, including framework wrappers such
as Lightning, without further dependencies, and either can be composed
with conventional augmentations (random crop, flip, colour jitter) applied
before or after the masking step.
